\begin{table}[htbp]
\centering
\small
\caption{Estatísticas Descritivas das Variáveis de Estudo (184 Municípios do Ceará)}
\label{tab:estatisticas_descritivas}
\begin{tabular}{llrrrrr}
\toprule
\textbf{Variável} & \textbf{Unidade} & \textbf{Média} & \textbf{Desv. Padrão} & \textbf{Mínimo} & \textbf{Mediana} & \textbf{Máximo} \\
\midrule
\multicolumn{7}{l}{\textit{\textbf{Desfechos Epidemiológicos}}} \\
Casos de COVID-19 por 100 mil hab. & casos por 100 mil habitantes & 1.794,6 & 1.231,0 & 86,21 & 1.575,9 & 7.582,6 \\
Óbitos por COVID-19 por 100 mil hab. & obitos por 100 mil habitantes & 48,29 & 32,56 & 0,00 & 42,83 & 138,20 \\
Taxa de letalidade (óbitos/casos) & proporcao (obitos/casos) & 0,03 & 0,02 & 0,00 & 0,03 & 0,17 \\
\midrule
\multicolumn{7}{l}{\textit{\textbf{Covariáveis Demográficas, Socioeconômicas e de Saúde}}} \\
População estimada (2019) & habitantes & 49.630,9 & 198.896,6 & 4.844,0 & 21.999,5 & 2.669.342,0 \\
Área territorial & km2 & 809,21 & 752,78 & 72,67 & 610,63 & 4.262,3 \\
Densidade demográfica & habitantes/km2 & 125,99 & 652,58 & 7,13 & 39,20 & 8.545,9 \\
PIB per capita & R$/habitante & 10.761,1 & 7.984,2 & 5.351,8 & 8.275,7 & 87.086,0 \\
IDHM (2010) & indice 0-1 & 0,62 & 0,03 & 0,54 & 0,61 & 0,75 \\
IDM (2018) & indice IPECE & 24,94 & 9,51 & 9,05 & 22,70 & 65,50 \\
Despesa municipal com saúde & R$ por 100 mil habitantes & 66.595.788,8 & 20.092.837,7 & 19.643.436,2 & 62.906.038,0 & 198.486.365,3 \\
Despesa municipal com educação & R$ por 100 mil habitantes & 96.593.085,7 & 25.036.453,7 & 37.592.999,4 & 91.393.494,5 & 210.694.228,9 \\
IDEB anos iniciais (fundamental) & indice IDEB & 6,37 & 0,86 & 4,70 & 6,20 & 9,40 \\
IDEB anos finais (fundamental) & indice IDEB & 5,24 & 0,60 & 3,90 & 5,10 & 7,80 \\
IDEB 3º ano (médio) & indice IDEB & 4,35 & 0,38 & 3,50 & 4,40 & 5,60 \\
Cobertura urbana de água & percentual & 97,76 & 2,73 & 87,33 & 98,74 & 100,00 \\
Consumo de energia elétrica & MWh por 100 mil habitantes & 109.623,5 & 64.090,8 & 23.356,5 & 87.100,3 & 451.049,3 \\
Taxa de homicídios & homicidios por 100 mil habitantes & 22,90 & 19,61 & 0,00 & 18,61 & 122,77 \\
Proporção de vereadoras eleitas & percentual & 31,85 & 2,66 & 26,67 & 31,58 & 48,15 \\
População no Bolsa Família & percentual & 46,05 & 8,91 & 22,56 & 47,28 & 64,03 \\
Unidades SUS por 1.000 hab. & unidades por mil habitantes & 0,72 & 0,27 & 0,12 & 0,68 & 1,71 \\
Leitos SUS por 1.000 hab. & leitos por mil habitantes & 1,44 & 0,95 & 0,00 & 1,24 & 6,11 \\
Profissionais SUS por 1.000 hab. & profissionais por mil habitantes & 9,16 & 2,16 & 5,29 & 8,80 & 20,56 \\
Região Metropolitana (binária) & indicador 0/1 & 0,10 & 0,31 & 0,00 & 0,00 & 1,00 \\
Semiárido (binária) & indicador 0/1 & 0,73 & 0,45 & 0,00 & 1,00 & 1,00 \\
População rural & proporcao 0-1 & 0,44 & 0,16 & 0,00 & 0,45 & 0,76 \\
Alinhamento partidário federal (binária) & indicador 0/1 & 0,51 & 0,50 & 0,00 & 1,00 & 1,00 \\
Votos válidos do prefeito eleito & proporcao 0-1 & 0,58 & 0,13 & 0,34 & 0,54 & 1,00 \\
IVS Infraestrutura Urbana & indice 0-1 & 0,28 & 0,11 & 0,06 & 0,27 & 0,62 \\
IVS Capital Humano & indice 0-1 & 0,51 & 0,06 & 0,33 & 0,51 & 0,64 \\
IVS Renda e Trabalho & indice 0-1 & 0,53 & 0,07 & 0,28 & 0,53 & 0,78 \\
População com emprego formal & proporcao 0-1 & 0,09 & 0,07 & 0,03 & 0,07 & 0,73 \\
População em extrema pobreza & proporcao 0-1 & 0,27 & 0,09 & 0,05 & 0,28 & 0,47 \\
Cobertura vacinal (< 1 ano) & proporcao (pode exceder 1; documentado como pendencia) & 0,86 & 0,13 & 0,20 & 0,89 & 1,07 \\
População com 60 anos ou mais & proporcao 0-1 & 0,14 & 0,03 & 0,05 & 0,13 & 0,22 \\
Internações circulatórias & internacoes por 100 mil habitantes & 411,39 & 167,27 & 135,51 & 367,82 & 1.115,4 \\
Internações respiratórias & internacoes por 100 mil habitantes & 657,34 & 385,43 & 136,00 & 546,38 & 2.754,9 \\
Internações por diabetes & internacoes por 100 mil habitantes & 68,16 & 47,12 & 6,11 & 57,92 & 284,65 \\
Internações por asma & internacoes por 100 mil habitantes & 27,74 & 56,52 & 0,00 & 13,55 & 640,87 \\
Transferências federais COVID & R$ declarados por 100 mil habitantes & 8.964.198,7 & 30.011.443,7 & 225.330,2 & 4.951.907,5 & 405.048.949,1 \\
Auxílio emergencial & R$ declarados por 100 mil habitantes & 116.565.459,6 & 173.048.020,7 & 3.128.724,2 & 97.607.590,2 & 2.294.432.962,7 \\
\bottomrule
\end{tabular}
\begin{tablenotes}\footnotesize
\item \textit{Fonte:} Elaboração própria a partir de dados do DATASUS, IBGE, IPECE, IPEA e INEP.
\end{tablenotes}
\end{table}